Current lie-detection methods for large language models (LLMs) treat all false outputs as a single category, yet the mechanisms behind these errors differ fundamentally in nature, detectability, and risk. We introduce a behavioral classification framework that distinguishes three types of LLM falsehoods: \emph{confabulations} arising from epistemic uncertainty, \emph{incentive-driven errors} where the model overrides its own knowledge under social pressure, and \emph{systematic errors} reflecting memorized misconceptions. Using multi-sample consistency analysis and paired pressure testing on 900 question-answer pairs across two state-of-the-art models (\gptfour and \gemini), we find that approximately 61\% of false outputs are confabulations, 18--23\% are incentive-driven, and 15--21\% are systematic errors. This distribution is remarkably consistent across both models ($\chi^2 = 1.23$, $p = 0.54$), suggesting it reflects fundamental properties of RLHF-trained models. A logistic regression classifier achieves 71.9\% cross-validated accuracy at distinguishing these three types from text features alone, confirming they produce partially but not fully separable behavioral signatures. Our findings reveal that roughly 1 in 5 LLM errors involves the model overriding its own knowledge---a qualitatively more concerning failure mode than simple ignorance, with direct implications for safety-critical deployment and detector design.
